\documentclass[10pt,a4paper]{amsart}
\usepackage{amsmath,amssymb,amsthm,graphicx}
\usepackage{verbatim}
\newenvironment{reference}{\comment}{\endcomment}
%\newenvironment{reference}{}{}
\newenvironment{slogan}{\comment}{\endcomment}
\newenvironment{history}{\comment}{\endcomment}

% For commutative diagrams we use Xy-pic
\usepackage[all]{xy}

% We use 2cell for 2-commutative diagrams.
\xyoption{2cell}
% Theorem environments.
%
\theoremstyle{plain}
\newtheorem{theorem}[subsection]{Theorem}
\newtheorem{proposition}[subsection]{Proposition}
\newtheorem{lemma}[subsection]{Lemma}

\theoremstyle{definition}
\newtheorem{definition}[subsection]{Definition}
\newtheorem{example}[subsection]{Example}
\newtheorem{exercise}[subsection]{Exercise}
\newtheorem{situation}[subsection]{Situation}

\theoremstyle{remark}
\newtheorem{remark}[subsection]{Remark}
\newtheorem{remarks}[subsection]{Remarks}

\numberwithin{equation}{subsection}

% Macros
%
\def\lim{\mathop{\mathrm{lim}}\nolimits}
\def\colim{\mathop{\mathrm{colim}}\nolimits}
\def\Spec{\mathop{\mathrm{Spec}}}
\def\Hom{\mathop{\mathrm{Hom}}\nolimits}
\def\Ext{\mathop{\mathrm{Ext}}\nolimits}
\def\SheafHom{\mathop{\mathcal{H}\!\mathit{om}}\nolimits}
\def\SheafExt{\mathop{\mathcal{E}\!\mathit{xt}}\nolimits}
\def\Sch{\mathit{Sch}}
\def\Mor{\mathop{\mathrm{Mor}}\nolimits}
\def\Ob{\mathop{\mathrm{Ob}}\nolimits}
\def\Sh{\mathop{\mathit{Sh}}\nolimits}
\def\NL{\mathop{N\!L}\nolimits}
\def\CH{\mathop{\mathrm{CH}}\nolimits}
\def\proetale{{pro\text{-}\acute{e}tale}}
\def\etale{{\acute{e}tale}}
\def\QCoh{\mathit{QCoh}}
\def\Ker{\mathop{\mathrm{Ker}}}
\def\Im{\mathop{\mathrm{Im}}}
\def\Coker{\mathop{\mathrm{Coker}}}
\def\Coim{\mathop{\mathrm{Coim}}}

% Boxtimes
%
\DeclareMathSymbol{\boxtimes}{\mathbin}{AMSa}{"02}

%
% Macros for moduli stacks/spaces
%
\def\QCohstack{\mathcal{QC}\!\mathit{oh}}
\def\Cohstack{\mathcal{C}\!\mathit{oh}}
\def\Spacesstack{\mathcal{S}\!\mathit{paces}}
\def\Quotfunctor{\mathrm{Quot}}
\def\Hilbfunctor{\mathrm{Hilb}}
\def\Curvesstack{\mathcal{C}\!\mathit{urves}}
\def\Polarizedstack{\mathcal{P}\!\mathit{olarized}}
\def\Complexesstack{\mathcal{C}\!\mathit{omplexes}}
% \Pic is the operator that assigns to X its picard group, usage \Pic(X)
% \Picardstack_{X/B} denotes the Picard stack of X over B
% \Picardfunctor_{X/B} denotes the Picard functor of X over B
\def\Pic{\mathop{\mathrm{Pic}}\nolimits}
\def\Picardstack{\mathcal{P}\!\mathit{ic}}
\def\Picardfunctor{\mathrm{Pic}}
\def\Deformationcategory{\mathcal{D}\!\mathit{ef}}
\usepackage{cleveref}
\crefname{originallemma}{original lemma}{original lemmas}
\Crefname{originallemma}{Original lemma}{Original lemmas}
\makeatletter
\newlabel{more-algebra-lemma-compose}{{0A5Y}{source}{}{sourcepointer-1}{}}
\newlabel{more-algebra-lemma-compose@cref}{{[originallemma][0][]0A5Y}{[1][0][]source}{}{}{}}
\makeatother
\crefname{originalsection}{original section}{original sections}
\Crefname{originalsection}{Original section}{Original sections}
\makeatletter
\newlabel{section-tensor-product}{{00CV}{source}{}{sourcepointer-2}{}}
\newlabel{section-tensor-product@cref}{{[originalsection][0][]00CV}{[1][0][]source}{}{}{}}
\makeatother
\title{Tensor-Hom adjunction for differential graded modules over a DGA}
\author{The Stacks Project Authors}
\date{stacks-project:dga.tex:local-snapshot-20261012; GNU FDL 1.2 or later}
\begin{document}
\maketitle
\noindent\textit{Editorial note.} All delivered mathematical text is the author's own, copied from the pinned Stacks Project chapter file sources/stacks-dga/pinned-original-source.tex; no statement or proof step was rewritten, reordered or removed.\par\medskip
\noindent\textit{Editorial note.} The counted proof is the author's complete argument: the right-hand identification is the defining universal property recorded in the same unit's tensor-product section (in-source reference); the left-hand unpacking uses the standard More on Algebra composition lemma for total complexes (navigation stub); both dependencies are stated in the proof and carried as stubs.\par\medskip
\noindent\textit{Editorial note.} The label lemma-characterize-hom names the counted statement, and the proof environment immediately following it in the source is delivered verbatim as its proof.\par\medskip
\noindent\textit{Editorial note.} The counted statement is carried by the author's own lemma environment in the chapter source.\par\medskip
\noindent\textit{Allowed standard prerequisites (editorial).}\begin{itemize}
\item Universal property of the tensor product. Same unit, Section section-tensor-product. The proof identifies the right hand side with differential-compatible A-bilinear maps via the defining property of the tensor product of differential graded modules recorded in that section.
\item Composition of total complexes. More on Algebra, Lemma more-algebra-lemma-compose (standard tooling). Left hand side: an element g of Hom(M', Hom(M, N)) determines a map of complexes g': Tot(M' tensor\_R M) -> N by the composition lemma for total complexes.
\end{itemize}\medskip
\ifdefined\hypertarget\hypertarget{sourcepointer-1}{}\fi
\noindent\textit{Source reference.} Original lemma 0A5Y: sources/stacks-chow/tags.txt, source lines 12826--12826. Prerequisite pointer invoked at proof line(s) 1782; it names a standard fact of a sibling chapter that the counted argument uses, and every step of the counted proof built on it is retained verbatim in the delivered body.\par\medskip
\ifdefined\hypertarget\hypertarget{sourcepointer-2}{}\fi
\noindent\textit{Source reference.} Original section 00CV: sources/stacks-dga/pinned-original-source.tex, source lines 1606--1606. Prerequisite pointer invoked at proof line(s) 1771; it names a standard fact of the same chapter that the counted argument uses, and every step of the counted proof built on it is retained verbatim in the delivered body.\par\medskip
\begin{lemma}
\label{lemma-characterize-hom}
Let $R$ be a ring. Let $(A, \text{d})$ be a differential graded $R$-algebra.
Let $M'$ be a right differential graded $A$-module and let
$M$ be a left differential graded $A$-module.
Let $N^\bullet$ be a complex of $R$-modules. Then we have
$$
\Hom_{\text{Mod}_{(A, d)}}(M', \Hom(M, N^\bullet)) =
\Hom_{\text{Comp}(R)}(M' \otimes_A M, N^\bullet)
$$
where $M \otimes_A M$ is viewed as a complex of $R$-modules
as in Section \ref{section-tensor-product}.
\end{lemma}

\begin{proof}
Let us show that both sides correspond to graded $A$-bilinear maps
$$
M' \times M \longrightarrow N^\bullet
$$
compatible with differentials. We have seen this is true for the right
hand side in Section \ref{section-tensor-product}. Given an element
$g$ of the left hand side, the equality of
More on Algebra, Lemma \ref{more-algebra-lemma-compose}
determines a map of complexes of $R$-modules
$g' : \text{Tot}(M' \otimes_R M) \to N^\bullet$.
In other words, we obtain a graded $R$-bilinear
map $g'' : M' \times M \to N^\bullet$ compatible with differentials.
The $A$-linearity of $g$ translates immediately
into $A$-bilinarity of $g''$.
\end{proof}

\end{document}
